\documentclass[10pt]{article}
\usepackage[a4paper,margin=24mm]{geometry}
\usepackage[T1]{fontenc}
\usepackage{amsmath,amssymb,booktabs,xcolor,hyperref}
\setlength{\parindent}{0pt}
\setlength{\parskip}{5pt}
\newcommand{\auditbox}[3]{\par\noindent\fcolorbox{#1}{#1!5}{\parbox{0.95\linewidth}{\textbf{#2}\par #3}}\par}
\title{Analytic Audit of Appendix A}
\author{DDLALM / SPLALM manuscript}
\date{13 October 2026}
\begin{document}
\maketitle
\auditbox{green!50!black}{PASS after one explanatory addition}{The parameter bounds,
problem-class membership, spectral decomposition, transient estimates, and
rotating-branch KKT lower bound are valid under the stated hypotheses.
The executor and an independent reviewer checked the derivations analytically.
No numerical verification was performed.}

\section*{Scope and finding}
This audit covers \texttt{appendix/LowerBoundRotation.tex} and its direct
interfaces in the lower-bound statements, problem-class definition, and rotating
trajectory proof. It is not a whole-manuscript audit.

\begin{center}
\begin{tabular}{llll}
\toprule
Finding & Severity / impact & Action & Status\\
\midrule
A-DENSITY-1 & Minor / exposition & Add derivation & Closed\\
\bottomrule
\end{tabular}
\end{center}

\auditbox{red!70!black}{Before}{The original text asserted that algebraic
exceptional sets have dense complement. It did not show that the repeated-root
condition is nontrivial or that boundary perturbations preserve admissibility.}
\auditbox{orange!80!black}{Why an addition was needed}{An algebraic equality can
hold identically. The actual parameter formulas must exclude that possibility.
The claimed density is correct; no counterexample or false constant was found.}
\auditbox{green!50!black}{After}{The proof now checks the relevant denominators,
exhibits a nonzero limit of the discriminant, and constructs admissible
perturbations. No assumption or theorem conclusion was changed.}

\section*{Complete added argument}
Use the appendix notation $D=\alpha W$, $V=\beta+d_\tau(1-\beta)$,
$\vartheta=(1-V)\xi$, and
\[
Q_{\rm f}(t)=(t-1+D)(t-1+V)-\vartheta(t-1+\alpha).
\]
For $\alpha,\beta,d_\tau\in(0,1)$ and $V\ne\alpha$, the definitions give
$W-1>0$ and $\vartheta>0$. Hence
\[
Q_{\rm f}(1-\alpha)=\alpha(W-1)(V-\alpha)\ne0,
\qquad Q_{\rm f}(1-V)=\vartheta(V-\alpha)\ne0.
\]
Also, the parameter definitions give
\[
\beta\eta\tau b^2=\nu\bigl(WV-(1-V)\xi\bigr),\qquad
Q_{\rm f}(1)=\frac{\alpha\beta\eta\tau b^2}{\nu}>0.
\]
These evaluations exclude all root-dependent denominators. The synchronization
system has determinant $\alpha(1-V)(\alpha-V)(W+\xi)\ne0$.

For fixed interior $\alpha,\beta,d_\tau$, the discriminant
\[
(D+V-\vartheta)^2-4(DV-\alpha\vartheta)
\]
is rational in $\epsilon^2$ and tends to $(U-V)^2$ as $\epsilon\downarrow0$.
For $U\ne V$ it is not identically zero and its zeros are isolated. Interior
variation of $\beta$ changes $V$ strictly, so $V=\alpha$ and $V=U$ can also
be avoided by arbitrarily small perturbations. If
$\epsilon(\rho+s)=c_1$, first decrease $\epsilon$ slightly to make room for
perturbing the other parameters. The strict inequalities
$\epsilon<c_2$ and $\alpha>\alpha_\epsilon$ are preserved for sufficiently
small perturbations. Thus admissible nondegenerate parameters approximate
every allowed endpoint and repeated-root case.

All entries of $Q_{\rm f}(\mathbb T)\mathbf r$ are continuous there because
$\alpha,b,V,W+\xi>0$ and the complex denominators are nonzero. The annihilator
identity therefore extends to every admissible tuple.

\section*{Other verified proof obligations}
\textbf{Parameter bounds.} The alpha threshold gives
$\nu/\alpha^2<1/64$ and $64U/65\le D\le U\le1$.
The first term in $c_1$ gives $\xi C_{\rm tr}\le1/16$ and hence
\[
\frac{12}{13}\le\frac{b^2}{64\epsilon^2}\le1,
\qquad 7\epsilon\le b\le8\epsilon.
\]
The real-part identity in the appendix yields both amplitude estimates.
The remaining terms in $c_1,c_2$ yield $\kappa\le1/32$ and the required
products involving $bs$ and $b\rho$.

\textbf{Spectral decomposition.} Eliminating the center coordinates gives the
stated quartic characteristic polynomial. Comparing coefficients using the
parameter identities gives its two quadratic factors. The synchronized
two-dimensional system yields exactly the stated $r_x,r_\lambda$.
For $|t|\ge1$,
\[
|(t-1+D)(t-1+V)|\ge DV|t|^2
>2\vartheta|t|\ge|\vartheta(t-1+\alpha)|.
\]
Thus every fast root is strictly inside the unit circle, including the
repeated-root cases. No diagonalizability assumption is used.

\textbf{Transient bounds.} The nonnegative convolution kernel has supremum
at most $1/D$ and sum $1/(DV)$. Finite-prefix maximum and sum inequalities
give
\[
|p_k|\le X/4,\qquad \sum_{k\ge0}|p_k|\le X/(4DV).
\]
The dual difference recursion and the weighted-average tail sum then give
$|\ell_k|\le s(\gamma/\eta+5/4)bX$. These bounds do not require a uniform
spectral gap as $\beta\downarrow0$.

\textbf{Membership and trajectory.} The radial pieces match through first
derivatives, and their Hessian norms are at most $\kappa$. The objective gap is
\[
F(0)-\inf F\le6\epsilon+32\gamma\epsilon^2\le\Delta.
\]
The trajectory induction verifies the current and staggered proximal conditions
in order, with both norms at most $3/4$. The two amplitude cases yield
\[
R_k\ge\frac b2\max\{X,|\Lambda|\}=b/6\ge7\epsilon/6.
\]

\section*{Closure and provenance}
Dependency order: parameter bounds, spectral identities, transient bounds,
then membership and trajectory. The forward reference in the membership proof
creates no cycle. Analytic endpoint checks covered $\alpha=1$, $\beta=1$,
$s=\tau$, $V=\alpha$, repeated roots, $\eta=\gamma$, and $\beta\downarrow0$.
No analytic counterexample was found. There are no open obligations in this scope.

The configured external reviewer model was unavailable. Independent review was
completed by the session agent \texttt{appendix\_a\_analytic\_review}, which also
verified the added paragraph. The manuscript compiles to 38 pages with no
undefined references or LaTeX errors in the final log. Detailed Chinese
derivations are in \texttt{notes/progress/appendix-a-audit-20261013.md}.
\end{document}
